\documentclass[10pt,a4paper]{amsart}
\usepackage[margin=19mm]{geometry}
\usepackage{amsmath,amssymb,mathtools}
\usepackage[scr=rsfs,cal=euler]{mathalfa}
\usepackage{xfrac}
\usepackage[unicode,hidelinks,bookmarks=false]{hyperref}
\newcommand{\ie}{i.e.\ }
\newcommand{\cf}{cf.\ }
\newcommand{\resp}{respectively\ }
\newcommand{\Subsection}{\subsection}
\providecommand{\nobreakdash}{\nobreak\hbox{-}}
\setlength{\emergencystretch}{2em}
\multlinegap0pt
\usepackage{amssymb}
\usepackage{bm}
\usepackage{algorithm}\usepackage{algpseudocode}
\usepackage{tikz}
\usepackage[shortlabels]{enumitem}
\usepackage{xcolor}
\newtheorem{lemma}{Lemma}
\title{Feasibility Determination for Subjective Probability Constraints}
\author{Taehoon Kim, Sigr\'un Andrad\'ottir, Seong-Hee Kim, Yuwei Zhou}
\date{Source: arXiv:2605.27237v2 (CC BY 4.0)}
\begin{document}
\maketitle

\noindent\emph{Source and licence.} Feasibility Determination for Subjective Probability Constraints. Version arXiv:2605.27237v2 (CC BY 4.0), distributed by arXiv under the Creative Commons Attribution 4.0 International licence, http://creativecommons.org/licenses/by/4.0/, as recorded in the arXiv licence metadata for this exact version and shown on the abstract page https://arxiv.org/abs/2605.27237v2. Full licence notice: \texttt{licenses/arxiv-2605.27237-CC-BY-4.0.txt}.\par\smallskip

\noindent\textit{Editorial scope.} Counted unit: the article's lemma thm:single\_system\_multiple\_threshold (source0.tex lines 663--672). The article's complete original proof of it is delivered with it (source0.tex lines 674--714, one \texttt{proof} environment of 41 lines). No mathematics is written, repaired or re-derived: the statement and the proof are the article's own lines, in the article's own notation, and no retained line is substituted or re-flowed.  Retained uncounted as the source-local context the proof needs: lines 596--599.  Nothing was omitted from any retained span, so the package carries no omission record.  No retained line contains a citation, so the wrapper carries no bibliography and no bibliography entry.  No retained line contains a figure environment, an image-inclusion command or drawing code, so no image is delivered and the package declares no asset dependency.  Editorial formatting only, in addition: an AMS \texttt{amsart} preamble that restates the article's own declarations and macros used by the retained lines, namely the environments \texttt{lemma} on separate counters, together with the standard packages those lines need.  The retained environments print in this document as Lemma 1 in order of appearance, whereas the article prints the same items with the numbering of its own full document.  The complete native source of the article as distributed in the arXiv e-print for this exact version is bundled unchanged as \texttt{sources/201477/source0.tex}.
\begin{lemma} 
\label{thm:single_system_single_threshold}
For a single system, a constraint $\ell$ with threshold $h$, $\mathcal{BRF}^{(1)}$ guarantees $\Pr\left( {\rm CD}_\ell(h_\ell) \right)  \ge 1 - \beta_\ell$.
\end{lemma}

\begin{lemma} 
\label{thm:single_system_multiple_threshold}
For a single system and a single constraint with thresholds $h_{\ell,1},h_{\ell,2},\ldots,h_{\ell, d_\ell}$, where $d_\ell\geq 1$, $\mathcal{BRF}^{(1)}$ guarantees 
\[
\Pr\left( {\rm CD}_\ell \right) = \Pr\left( \cap_{m=1}^{d_\ell}{\rm CD}_\ell(h_{\ell, m}) \right)  \ge \left\{ \begin{array}{ll} 
1 - \beta_\ell, & \mbox{if } d_\ell=1; \\
1 - 2\beta_\ell, & \mbox{if } d_\ell\ge 2.
\end{array} \right.
\]
\end{lemma}

\begin{proof}
    \renewcommand{\qedsymbol}{}
    If $d_\ell=1$, then 
\[
\Pr\left( {\rm CD}_\ell \right) = \Pr\left( \cap_{m=1}^{d_\ell}{\rm CD}_{\ell}(h_{\ell, m}) \right)  = \Pr\left( {\rm CD}(h_{\ell, 1}) \right) \ge 1 - \beta_\ell,
\]
where the inequality holds due to Lemma~\ref{thm:single_system_single_threshold}.

If $d_\ell \ge 2$, there are three possible cases regarding the feasibility of the system. 
\begin{itemize}
    \item[] {\bf Case 1:}  The system is desirable with respect to all thresholds. Since $h_{\ell,1} < h_{\ell,2} < \cdots < h_{\ell,d_\ell}$ and we use the same random number $U_{n}$ to generate the Bernoulli random variables $I_{\ell mn}$ for $m=1,2,\ldots,d_\ell$, we have $I_{\ell 1n} \le I_{\ell 2n} \le \cdots \le I_{\ell d_\ell n}$, for $n=1,2,\ldots$. It follows that
    \begin{equation} \label{eqn:dummy}
    \sum_{n=1}^{r}\left(Y_{\ell n} -I_{\ell 1n}\right) \ge \sum_{n=1}^{r}\left(Y_{\ell n} -I_{\ell 2n}\right)\ge \cdots \ge \sum_{n=1}^{r}\left(Y_{\ell n} -I_{\ell d n}\right).
    \end{equation}
    A correct decision with respect to $h_{\ell, m}$ is made if $\sum_{n=1}^{r}(Y_{\ell n} -I_{\ell mn})$ reaches $-H_\ell$ before $H_\ell$. Equation \eqref{eqn:dummy} implies that if we have a correct decision at $h_{\ell, 1}$, i.e.,  if $\sum_{n=1}^{r}(Y_{\ell n} -I_{\ell 1n})$ reaches $-H_\ell$ before $H_\ell$, then we also have correct decisions for the other thresholds $h_{\ell, 2}< \cdots < h_{\ell, d_\ell}$. That is, we have ${\rm CD}(h_{\ell,1}) \subseteq {\rm CD}(h_{\ell, 2}) \subseteq \cdots \subseteq {\rm CD}(h_{\ell, d_\ell})$. Therefore, Lemma~\ref{thm:single_system_single_threshold} yields
    \[ 
    \Pr\left( {\rm CD}_\ell \right)=\Pr\Big(\cap_{m=1}^{d_\ell} {\rm CD}_\ell(h_{\ell, m})\Big) = \Pr\Big( {\rm CD}_\ell(h_{\ell,1})\Big) \ge 1 - \beta_\ell \ge 1 - 2 \beta_\ell.
    \] 
    \item[] {\bf Case 2:} The system is unacceptable with respect to all thresholds. A correct decision occurs when $\sum_{n=1}^{r}(Y_{\ell n} -I_{\ell mn})$ reaches $H_\ell$ before $-H_\ell$. Similarly to {\bf Case 1}, if a correct decision is made for $h_{\ell, d_\ell}$, correct decisions are also made for thresholds $h_{\ell, 1}\leq \cdots\leq h_{\ell, d_\ell}$, i.e., ${\rm CD}_\ell(h_{\ell, d_\ell}) \subseteq {\rm CD}_\ell(h_{\ell,d_\ell-1}) \subseteq \cdots \subseteq {\rm CD}_\ell(h_{\ell,1})$. Thus, it follows from Lemma~\ref{thm:single_system_single_threshold} that
    \[ 
    \Pr\left( {\rm CD}_\ell \right)=\Pr\Big(\cap_{m=1}^{d_\ell} {\rm CD}_\ell(h_{\ell,m})\Big) = \Pr\Big( {\rm CD}_\ell(h_{\ell,d_\ell})\Big) \ge 1 - \beta_\ell \ge 1 - 2\beta_\ell.
    \]
    \item[] {\bf Case 3:} The system is unacceptable with respect to thresholds $h_{\ell, 1} < \cdots < h_{\ell, \underline{\kappa}}$ and is desirable with respect to thresholds $h_{\ell, \overline{\kappa}} < \cdots < h_{\ell,d_\ell}$, where $1\leq \underline{\kappa} <  \overline{\kappa}\leq d_\ell$.
Suppose that we add two additional thresholds $h_\ell^L$ and $h_\ell^U$ such that 
\begin{equation} 
\label{eqn:LBUB}
\frac{h_\ell^U (1-p_\ell)}{(1-h_\ell^U)p_\ell}  = \theta_\ell \quad \mbox{and} \quad {(1-h_\ell^L)p_\ell \over h_\ell^L(1-p_\ell)} =\theta_\ell,
\end{equation}
i.e., these thresholds have odds-ratios with respect to $p_\ell$ that exactly equal $\theta_\ell$. Then we have that the system is desirable for $h_\ell^U$, unacceptable for $h_\ell^L$, and 
\[
h_{\ell, 1} < \cdots < h_{\ell, \underline{\kappa}} \le h_\ell^L < h_{\ell,\underline{\kappa}+1} < \cdots < h_{\ell,\overline{\kappa}-1} < h_\ell^U \le h_{\ell,\overline{\kappa}} < \cdots < h_{\ell,d_\ell}.
\]
A correct decision occurs when the system is deemed infeasible with respect to $h_{\ell,1}, \ldots, h_{\ell,\underline{\kappa}}$ but feasible with respect to $h_{\ell,\overline{\kappa}}, \ldots, h_{\ell,d_\ell}$. Consider dummy Bernoulli random variables $I_{\ell n}^{L}$ and $I_{\ell n}^{U}$ that are generated based on thresholds $h_\ell^{L}$ and $h_\ell^{U}$, respectively, applying the same random number $U_{n}$ used for generating $I_{\ell 1n},\ldots, I_{\ell d_\ell n}$, which yields $I_{\ell 1n}\leq \cdots\leq I_{\ell \underline{\kappa} n}\leq I_{\ell n}^L\leq I_{\ell (\underline{\kappa}+1) n}\leq \cdots \leq I_{\ell (\overline{\kappa}-1) n}\leq I_{\ell n}^U \leq I_{\ell \overline{\kappa} n}\leq \cdots \leq I_{\ell d_\ell n}$ for $n=1,2,\ldots$.
It follows that ${\rm CD}_\ell(h_\ell^{L}) \subseteq {\rm CD}_\ell(h_{\ell, \underline{\kappa}})  \subseteq \cdots  \subseteq {\rm CD}_\ell(h_{\ell,1})$ and ${\rm CD}_\ell(h_\ell^{U}) \subseteq {\rm CD}_\ell(h_{\ell,\overline{\kappa}}) \subseteq \cdots \subseteq  {\rm CD}_{\ell}(h_{\ell,d_\ell})$. For thresholds $h_{\ell,\underline{\kappa}+1}, \ldots, h_{ \ell, \overline{\kappa}-1}$, the system is acceptable, meaning that any feasibility decision is a correct decision. Therefore, the Bonferroni inequality and Lemma \ref{thm:single_system_single_threshold} yield
\begin{align*}
\Pr\left( {\rm CD}_\ell \right)  &= \Pr\Big(\cap_{m=1}^{d_\ell} {\rm CD}_\ell(h_{\ell, m})\Big) 
\geq \Pr\Big( {\rm CD}_\ell(h_\ell^{L}) \cap {\rm CD}_\ell(h_\ell^{U})\Big) \\
&\ge \Pr\Big( {\rm CD}_\ell(h_\ell^{L})\Big) + \Pr\Big({\rm CD}_\ell(h_\ell^{U})\Big) -1 \ge 1 - 2\beta_\ell. \tag*{$\square$}
\end{align*}
\end{itemize}
\end{proof}

\end{document}
